% Q02_flawed_answers.tex — GRADER CALIBRATION FIXTURE (deliberately flawed).
% A drop-in replacement for paper/Q02_answers.tex (spec §4.7, §6.12). B1 and B3 carry one
% planted flaw each; B2 is the clean control. Grade it with the rubric's probes
% (paper/rubrics/Q02_rubric.md), never by eye. It holds worked answers to Q02's written
% problems: open it only after Q02 is done.
%
% EXPECTED[B1][induction-soundness][2]: in (a)'s zero-top case the answer proves ToNat(p) = ToNat(s) by unfolding both at the LOW limb and claiming p[1..] = s[1..] for p = s[..|s| - 1]; that sequence equation is false. The finding must locate that step and exhibit s = [5, 0]: p = [5], p[1..] = [] while s[1..] = [0] (Probe A'': an @Test expecting s[..|s| - 1][1..] == s[1..] at [5, 0] prints FAILED; captured on the pins). The conclusion ToNat(p) = ToNat(s) is true; the justification is what fails. Statement-fidelity, case-exhaustiveness and gap-accounting are correct and must not score <= 2; blueprint-quality may score 2 only by exhibiting the same instance; (b) is correct and must draw no flaw.
% EXPECTED[B2][correctness][4]: correct, no finding (Probe B on its state: 1 passed, 0 failed; Probe B' on its loop with the companion clauses: verified; both captured on the pins). A grader that reports a flaw on B2 is hallucinating.
% EXPECTED[B3][card-fidelity][2]: Budget — the field says the loop retries without a fixed limit until Dafny accepts or a 30-minute wall clock runs out; pinned line 772 gives a budget of N = 10 attempts, and line 787 N = 5 for the held-out runs. The finding must quote line 772.
% EXPECTED[B3][card-fidelity][4]: Task shape, Oracle, What is trusted, Cheat surface, Metric, Variance and Artifact — correct, no finding (lines 770, 772, 775, 778, 780-787, 789). A grader that reports a flaw on any of them is hallucinating.
% Expected per-problem grades (min across dimensions or fields): B1 = 2, B2 = 4, B3 = 2.
% A calibration run FAILS if B1 or B3 scores >= 3, if a planted finding does not locate
% its step or field (with the instance or the quoted line named above), or if ANY flaw
% is reported on B2 or on a B3 field other than Budget.
\documentclass[11pt]{article}
\input{preamble}

\title{Q02 --- First syntax I \\ \large Written answers}
\author{learner}
\date{}

\begin{document}
\maketitle

% ---------------------------------------------------------------------------
\section*{Problem B1 (informal-proof-first)}\label{ans:B1}
\begin{answer}
\textbf{(a) Statement.} For every limb sequence $s$,
$\mathit{ToNat}(\mathit{Normalize}(s)) = \mathit{ToNat}(s)$. No hypothesis on
$s$.

\textbf{Induction on / decreases measure:} the length $|s|$, strong induction;
the lemma's \texttt{decreases |s|}.

\textbf{Base:} $|s| = 0$. Then $\mathit{Normalize}(s) = [\,]$, which is $s$,
so the two sides are the same number (0).

\textbf{IH (stated as a formula):} for every limb sequence $t$ with
$|t| < |s|$, $\mathit{ToNat}(\mathit{Normalize}(t)) = \mathit{ToNat}(t)$.

\textbf{Step} (one paragraph per case of \texttt{Normalize}):

\emph{Top limb non-zero} ($|s| > 0$, $s[|s|-1] \neq 0$):
$\mathit{Normalize}(s) = s$ by the second branch, so there is nothing to prove.

\emph{Top limb zero} ($|s| > 0$, $s[|s|-1] = 0$): let $p = s[..|s|-1]$, so
$|p| = |s| - 1$. The third branch gives $\mathit{Normalize}(s) =
\mathit{Normalize}(p)$, and the IH at $p$ gives
$\mathit{ToNat}(\mathit{Normalize}(p)) = \mathit{ToNat}(p)$. It remains to show
$\mathit{ToNat}(p) = \mathit{ToNat}(s)$. Unfold \texttt{ToNat} once on each
side: $\mathit{ToNat}(s) = s[0] + \mathit{BASE} \cdot \mathit{ToNat}(s[1..])$
and $\mathit{ToNat}(p) = p[0] + \mathit{BASE} \cdot \mathit{ToNat}(p[1..])$.
Here $p[0] = s[0]$, and $p[1..] = s[1..]$ because $p$ and $s$ agree on every
index below $|s| - 1$; so the two unfoldings are equal, and
$\mathit{ToNat}(\mathit{Normalize}(s)) = \mathit{ToNat}(p) = \mathit{ToNat}(s)$.

\textbf{Helper lemmas needed (this is the blueprint;} each by its statement,
and where it comes from\textbf{):} none beyond one unfolding of \texttt{ToNat}
on each side and the recursive call at $p$.

\textbf{Gaps I know about} (declared gaps score better than silent ones):
none.

\textbf{(b) Loop invariant (stated as a formula over \texttt{a} and
\texttt{old(a)}):} the number the limbs of \texttt{a} denote never changes:
$\mathit{seq\_to\_nat}(a@) = \mathit{seq\_to\_nat}(\mathit{old}(a)@)$ at the
head of every pass.

\textbf{What keeps it across one pass:} the guard says the vector is non-empty
and its top limb is zero; the given call states that dropping that limb keeps
\texttt{seq\_to\_nat}, and \texttt{pop} drops exactly that limb, so the value
after the pass is the value before it, hence still the value at the call.

\textbf{What gives each \texttt{ensures} clause at the exit:} the guard is
false, so the vector is empty or ends in a non-zero limb, which is
\texttt{wf()}; the invariant is the value clause. The loop stops because each
pass removes one limb (the measure \texttt{a.limbs.len()}).
\end{answer}

% ---------------------------------------------------------------------------
\section*{Problem B2 (concept)}\label{ans:B2}
\begin{answer}
\textbf{(a)} Inductiveness is decided by the first two obligations: the clause
must hold before the loop starts, and every pass that begins with it (and the
guard) true must end with it true. The third obligation, what the invariant
gives after the loop, says whether the invariant is strong enough for the
postcondition; a clause can be inductive and still too weak there.

\textbf{(b)} The loop: \texttt{k, t := 0, 0; while k < m \{ k, t := k + 1, t + 2; \}}
with \texttt{m: nat}. \\
The clause: $t \leq 2m$. \\
Why it holds in every reachable state: every reachable state has $t = 2k$ and
$k \leq m$, so $t \leq 2m$. \\
The state before one pass, and the clause after it: $k = 0$, $t = 2$, $m = 1$.
The clause holds ($2 \leq 2$) and so does the guard ($0 < 1$); after the pass
$t = 4 > 2 = 2m$.

\textbf{(c)} With $t \leq 2m$ as the only clause, pruning finds it is not
preserved and drops it, true or not. Next to $k \leq m$ and $t = 2k$ it is
preserved, so pruning keeps it along with the other two.
\end{answer}

% ---------------------------------------------------------------------------
\section*{Problem B3 (paper card)}\label{ans:B3}
\begin{answer}
\textbf{Source:} cheatsheet entry 02.7, lines 766--791 (the pinned range).

\textbf{Task shape:} a Dafny program arrives with its pre- and postconditions
but no proof annotations; the model must add invariants, assertions and
termination measures, leaving the code as it is, until Dafny accepts it.

\textbf{Oracle:} the Dafny verifier accepting the program, plus a syntactic
check that removing every annotation from the attempt gives back the original
program.

\textbf{What is trusted:} Dafny and its SMT solver, the parser used to strip
annotations, and the benchmark specifications themselves; the card does not
say whether edited specifications or assumptions are rejected.

\textbf{Cheat surface:} rewriting the code is caught by the strip-and-compare
check; changing the specification or adding an assumption is not known to be
caught, and the hints may carry benchmark knowledge.

\textbf{Metric:} percentage of programs verified on Clover, MBPP, HumanEval and
DafnyBench, with 86.18\% on DafnyBench for Claude 3.5 as the headline, plus
ablations and a small-model result on held-out programs.

\textbf{Budget:} the loop retries without a fixed limit until Dafny accepts or a
30-minute wall clock runs out; no cost figures are reported.

\textbf{Variance:} a single run per configuration with two closed models, no
intervals; some numbers disagree between the text and the figures.

\textbf{Artifact:} no code is released; the distilled training set is
described but not stated to be released.
\end{answer}

\end{document}
